BepiColombo és una missió espacial que té per objectiu l'exploració de Mercuri. La missió va ser llançada l'any 2018, i hi arribarà el 2025. Un cop allà, posarà en òrbita dos satèl·lits al voltant del planeta. Un dels satèl·lits és el \textit{Mercury Planetary Orbiter} (MPO), construït per l'Agència Espacial Europea, que orbitarà al voltant de Mercuri amb un radi orbital mitjà de 3\,360~km.

\begin{apartat}{1,25}
Considereu un satèl·lit que fa una òrbita circular al voltant de Mercuri. Deduïu l'expressió de la velocitat orbital del satèl·lit en funció del radi orbital i la massa de Mercuri (indiqueu clarament en quins principis o lleis físiques us baseu per fer la vostra deducció). Amb aquesta expressió, calculeu la velocitat orbital del satèl·lit MPO mentre orbita al voltant de Mercuri. Calculeu quantes voltes haurà fet al planeta al cap d'un any terrestre.
\end{apartat}

\begin{apartat}{1,25}
A partir de l'expressió general de l'energia mecànica, obtingueu la seva equació per al cas particular d'un satèl·lit en òrbita circular (cal que l'equació final només estigui expressada en funció de $G$, el radi orbital i les masses del satèl·lit i del planeta). Una vegada el satèl·lit MPO estigui orbitant al voltant de Mercuri, encara tindrà combustible per a poder fer maniobres. Considereu que el combustible disponible pot proporcionar una energia de $4,5\times 10^{9}\ \mathrm{J}$. Determineu el valor màxim que podria tenir la massa del MPO per tal que amb l'energia disponible pogués escapar del camp gravitatori de Mercuri.
\end{apartat}

\dades{%
  $G = 6,67\times 10^{-11}\ \mathrm{N\,m^2\,kg^{-2}}$.\\
  Massa de Mercuri, $M_\mathrm{M} = 3,285\times 10^{23}\ \mathrm{kg}$.\\
  Any terrestre $= 365,25$ dies.}
